
\newcommand{\numTD}{TD5}
\newcommand{\themeTD}{Complexité}

\begin{center}
\begin{tabular}{|p{2cm}c|}
\hline
{\includegraphics[width=1.8cm,viewport=0 0 337 248]{../images/sorbonne.png}} & \raisebox{2ex}{\begin{Large}\textbf{M1SOL020}\end{Large}}\\
2019-2020& \raisebox{2ex}{\begin{Large}\textbf{ Epistémologie de l'Informatique}\end{Large}}\\
&  \begin{large}\textbf{\numTD}\end{large} \\
&  \begin{large} \textbf{\themeTD}\end{large} \\
& Gaël Lejeune, Sorbonne Université \\
& \tiny{Inspiré d'Agnès Delaborde 2015-2016}\\
\hline
\end{tabular}
\end{center}


\hrule
%%%%%%%%%%%%%%%%%%%%%%%%%EN-TETE%%%%%%%%%%%%%%%%%%%%%%%%%%%%%
%\renewcommand{\contentsname}{Sommaire du TD}
%\tableofcontents

\noindent\fcolorbox{red}{lightgray}{
\begin{minipage}{12cm}
\section*{Objectifs}

\begin{itemize}
  \item Comprendre la notion de complexité
  \item En déduire l'importance de l'algorithme
  \item Parcourir un tableau de façon efficace
\end{itemize}
\end{minipage}
}

\section{Principe}

\begin{figure}[h]
\includegraphics[width=.7\textwidth]{bigocheatsheet.png}
\caption{Principales classes de complexité (source : \url{bigocheatsheet.com})}
\end{figure}

\begin{table}[h]
\begin{tabular}{l|l|l}
Temps		&Dénomination		&Cas pratique\\
\hline
O(1)		& complexité constante&Accès à une cellule d'une matrice\\
O($\log$(n))	&complexité logarithmique& Recherche dichotomique\\
O(n)		&complexité linéaire&parcours d'une liste\\
O($n^{{2}}$)	&complexité quadratique (polynomiale)&Parcours d'une matrice, Distance de Levenshtein\\
2$^{{{\rm {{poly}(n)}}}}$complexité exponentielle& Désambiguïsation\\
\end{tabular}
\caption{Cas de complexité d'algorithmes (inspiré de Wikipedia \footnote{\url{https://fr.wikipedia.org/wiki/Analyse_de_la_complexit\%C3\%A9\_des\_algorithmes}})}
\end{table}
%%TODO: POur exponentielle complexité CCG

\section{Itérations}

\begin{itemize}
  \item  Rédiger l'algorithme (la séquence d'instructions) d'une boucle de 1 à 100.  A chaque pas de la boucle, une variable n prend la valeur du compteur. En sortie de boucle, on affiche n.
  \item  Réaliser le même algorithme pour une boucle de 1 à 100*1000
  \item  Codez ces deux algorithmes en \textsc{Python}
  \item  Mesurez la différence dans le temps d'exécution entre ces deux scripts : utilisez la méthode \texttt{clock} (version de Python avant 3.7) ou perf\_counter (seule disponible dans Python 3.7 et suivants) de la librairie \texttt{time} en stockant la valeur de \texttt{time.clock()} ou  \texttt{time.perf\_counter()}avant la boucle et en le comparant à sa valeur après la boucle.
  \item  Quelle est la complexité de l'algorithme (unité: nombre d'itérations) ?
\end{itemize}

\section{Parcours de tableau}

\begin{itemize}
  \item  Ecrivez un algorithme de création d'un tableau de 4 par 4. Chaque cellule contient une lettre, et chaque colonne représente l'alphabet (on s'arrête ici au bout de la 4ème lettre puisqu'il s'agit d'un tableau de 4 par 4).

\begin{tabular}{c|c|c|c}
A&A&A&A\\
\hline
B&B&B&B\\
\hline
C&C&C&C\\
\hline
D&D&D&D\\
\end{tabular}

  \item  Ecrivez un algorithme de parcours du tableau pour compter toutes les occurrences des lettres "A" et "C". Cet algorithme ne doit parcourir qu'une seule fois le tableau.
  \item  Rédigez le même algorithme, sauf que vous parcourrez le tableau une fois pour compter la présence de "A", et une seconde fois pour compter la présence des "C".
  \item  Quelle est la complexité de chacun des deux algorithmes (unité: nombre de passage dans la boucle) ?
  \item  Remarquez les temps d'exécution pour le(s) parcours du tableau. 
  \item  Si nous avions, au lieu d'un tableau 4 par 4, un tableau de 1000 par 2000, quelle serait la complexité des deux parcours ?
\end{itemize}

\section{Analyses de complexité}

 Quelle est la différence de complexité entre ces deux algorithmes qui vérifient la présence d'un entier dans un tableau (unité: nombre de passages dans la boucle tantque) ?

\vspace{0.5cm}

\textbf{Algorithme 1}
\vspace{-0.3cm}
\begin{verbatim}
e = 57
T = [x for x in range(100)]
i = 0
trouve =False
while trouve ==False and i <= len(T):
    if e == T[i]:
        trouve = True
    i += 1
\end{verbatim}


\textbf{Algorithme 2}
\vspace{-0.3cm}

\begin{verbatim}
e = 57
T = [x for x in range(100)]
i = 0
trouve =False
while i <= len(T):
    if e == T[i]:
        trouve = True
        break
    i +=1
\end{verbatim}
\vspace{-0.5cm}
\begin{itemize}
  \item  Indiquez, pour l'algorithme 1 ci-dessus, la complexité dans le meilleur des cas, dans le pire des cas, et la complexité moyenne.
  \item  Quelle est la complexité de l'algorithme ci-dessous ? (unité: passage dans la boucle)

\begin{python}
i, n, v = 0, 4, -1
while i<=n:
    v = v*2
    print(v)
\end{python}
  \item  Quelle est la complexité de l'algorithme ci-dessous ? (nombre d'opérations effectuées)
\begin{python}
i = 50
y = 4/3 * 4 * i
print(y)
\end{python}

  \item  Peut-on rendre plus efficace l'algorithme ci-dessous ?
\begin{verbatim}
def affichtab(tab)
    i, n = 0, len(tab)
    while i <= n:
        print(tab[i])
        i += 1

import string

s = "tableau d'entiers"
t = range(1, 11)
print("Voici le contenu de ", s)
affichtab(t)
s = "chaîne de caractères"

t = string.ascii_lowercase
print("Voici le contenu de ", s)
affichtab(t)
\end{verbatim}
\end{itemize}
%f. Proposez un algorithme efficace pour parcourir le tableau ci-dessous. Le parcours doit s'arrêter lorsqu'on a trouvé la combinaison [3][4][2].
%\begin{python}
%t[entiers de 0 à 4][entiers de 0 à 7][entiers de 0 à 2]

%\end{python}
